%% Preamble
% ---
\documentclass{article}
\usepackage[T1]{fontenc}
\usepackage[utf8]{inputenc}
\usepackage{array}
\usepackage{graphicx}
\newcommand{\safeincludegraphics}[2][]{\IfFileExists{#2}{\includegraphics[#1]{#2}}{\mbox{}}}
\graphicspath{{./}}
\usepackage{xcolor}
\usepackage[colorlinks]{hyperref}
\usepackage[margin=1in,left=1in,bottom=1in]{geometry}
\usepackage{multicol}
\usepackage{fancyhdr}
\usepackage{setspace}
\usepackage{indentfirst}
\hypersetup{
  colorlinks = true,
  linkcolor  = black
}
\usepackage{tocloft}


\begin{document}

	\begin{center}
		\thispagestyle{empty}
			\safeincludegraphics[width=3cm]{wce_logo.png}
			\\[0.20in]
			\huge \textbf{Walchand College Of Engineering, Sangli.}	\\[0.05in]
\Large  \textbf{(An Autonomous Institute)}
			\\[0.35in]
\textbf{Department Of\\ Computer Science and Engineering}
			%\\[0.30in]\Large \textbf{Year:2026-2027}
			\\[0.30in]{A Project Report \\ On}\\[0.30in]
\LARGE \textbf	{\color{blue}IdleGrid: Distributed Idle/Co-Tenant Compute\\Harvesting Platform for Campus Labs} 
			\\ [0.35in]
			\large	Submitted by\\[0.30in]
\begin{multicols}{2}	
			\textbf{Avadhut Mali\\ (PRN 23510003)}
			\end{multicols}
			\begin{multicols}{2}
		    \textbf{Yash Savalkar\\ (PRN 23510027)}
			\end{multicols} 
			\begin{multicols}{2}
		    \textbf{Krishna Dhumal\\ (PRN 23510057)}
            \end{multicols} 
			\begin{multicols}{2}
		    \textbf{Siddharth Swami\\ (PRN 23510096)}
            
			
			\end{multicols} 
			
			 Under the Guidance \\of \\	
			
			\begin{center}
			\large \textbf{[Guide name and affiliation]}
			\\[0.30in]\Large \textbf{Year: 2026 - 2027}
			\end{center}
			%\begin{multicols}
			%\\ \large \textbf{Dr.abc} \\HOD\\Computer Science and Engineering Dept,\\WCE, Sangli.	
			%\end{multicols}
		
			
%\vspace*{\fill}
			% \Large \textbf{May 2017}
	 \end{center}
	\pagebreak
	%\pagestyle{fancy}
	\fancyhf{}
	%\lhead{Time Series Data Analysis for Stock Market Prediction using Data Mining Techniques with R}
	\rfoot{\thepage}
\pagestyle{fancy}
	
	\begin{center}
		\safeincludegraphics[height=4cm]{wce_logo.png}
		\\
		\textbf{\Large Walchand College of Engineering, Sangli}\\
	\Large  (An Autonomous Institute)
\\[0.2in]
\textbf{Department\\ Of\\ Computer Science and Engineering}\\[0.4in]
\LARGE \textbf	{\color{blue}CERTIFICATE } \\[0.4in]
\end{center}


	\par
	\setlength{\parindent}{4em}
	%\doublespacing
	\space
This is to certify that the Project Report entitled, 
\textbf {"IdleGrid: Distributed Idle/Co-Tenant Compute Harvesting Platform for Campus Labs"} submitted by\\
\begin{multicols}{2}
\textbf{Avadhut Mali\\ (PRN 23510003)}
\end{multicols}
\begin{multicols}{2}
\textbf{Yash Savalkar\\ (PRN 23510027)}
\end{multicols}
\begin{multicols}{2}
\textbf{Krishna Dhumal\\ (PRN 23510057)}
\end{multicols}
\begin{multicols}{2}
\textbf{Siddharth Swami\\ (PRN 23510096)}
\end{multicols}
 to Walchand College of Engineering, Sangli, is a record of bona fide Project work carried out by them under my supervision and guidance and is worthy of consideration for the award of the degree of Bachelor of Technology in Computer Science and Engineering of the Institute. 
	

				\vspace{2cm}
\begin{multicols}{2}
	\begin{center}
	\textbf{[Guide name]}
	\\ Guide
	\end{center}
	\columnbreak
	\begin{center}
	External Examiner
	\end{center}
\end{multicols}

\vspace{1.5cm}
\begin{center}
	\textbf{Dr. A. R. Surve}
	\\ Head
	\\ Department of Computer Science and Engineering
\end{center}

	\pagebreak
	
	
	\pagebreak
		\pagebreak
	\begin{center}
	\Huge \textbf {\color {blue}Acknowledgement} \\
	\end{center}
	\vspace{1cm}
	{
	\onehalfspacing
	\setlength{\parindent}{2em}
	\setlength{\parskip}{0.8em}

We feel immense pleasure in submitting this project report entitled \textbf{``IdleGrid: Distributed Idle/Co-Tenant Compute Harvesting Platform for Campus Labs''}.

We are thankful to our guide, \textbf{[Guide name]}, for the valuable guidance and kind help during the completion of the project.

We are also thankful to the Head of the Department of Computer Science and Engineering, \textbf{Dr. A. R. Surve}, for the valuable guidance during the completion of the project.

We would like to thank all the faculty members and staff of the Department of Computer Science and Engineering for their generous help in various ways in the completion of this project.

We would like to thank all our friends, and especially our classmates, for the thoughtful and mind-stimulating discussions we had, which prompted us to think beyond the obvious. We have enjoyed their companionship so much during our stay at Walchand College of Engineering, Sangli.

	}
	
	\pagebreak
	\begin{center}
		\Huge \textbf{Declaration} \\
	\end{center}
	\vspace{1cm}
	\par 
	\setlength{\parindent}{8em}
	\doublespacing
	\doublespacing
    \noindent
We hereby declare that work presented in this project report titled 
\textbf {"IdleGrid: Distributed Idle/Co-Tenant Compute Harvesting Platform for Campus Labs"}
 submitted by us in the partial fulfillment of the requirement of the award of the degree of \textbf{Bachelor of Technology (B.Tech)} submitted in the \textbf{Department of Computer Science and Engineering, Walchand College of Engineering, Sangli}, is an authentic record of our project work carried out under the guidance of [Guide name and affiliation].
	

		\begin{multicols}{2}
\textbf{Avadhut Mali\\ (PRN 23510003)}
\end{multicols}
\begin{multicols}{2}
\textbf{Yash Savalkar\\ (PRN 23510027)}
\end{multicols}
\begin{multicols}{2}
\textbf{Krishna Dhumal\\ (PRN 23510057)}
\end{multicols}
\begin{multicols}{2}
\textbf{Siddharth Swami\\ (PRN 23510096)}
\end{multicols}

		\begin{multicols}{2}
		\noindent
		  Date : \\
		  Place : Sangli
		  \flushright{\vspace{1cm}(Signatures)}
		\end{multicols}
	
	\pagebreak
	\begin{center}
	\Large\textbf{ABSTRACT}
	\end{center}
	\noindent
Campus computer laboratories contain desktop computers whose capacity is often underused, while students may lack hardware for compute-intensive coursework. This project presents IdleGrid, a prototype platform that coordinates spare resources across campus PCs through a central Master service and per-PC Agents. A browser dashboard communicates with the Spring Boot Master, which maintains node and job state, accepts resource-constrained jobs, and assigns work using a first-fit scheduler. Agents report CPU, memory, and disk capacity, poll for assignments, and execute jobs in Docker containers managed through Docker Desktop and WSL2. The system also provides a shared-file workflow in which the Master stores primary file copies and eligible Agents hold best-effort replicas. A separate distributed-training workflow reserves two or three workers and runs a user-supplied PyTorch DistributedDataParallel (DDP) program using the CPU/Gloo backend. The implementation uses Java 17 and Spring Boot for the backend and Agent, and React, Vite, and Tailwind CSS for the dashboard. Node, job, booking, and training-run control records are held in memory; file data and its index are stored on the Master filesystem. The prototype is intended for a trusted campus network and does not include user authentication, durable job state, or checkpoint-based recovery. Its scheduling, resource isolation, file replication, and multi-node training capabilities should be evaluated on the target lab PCs before performance conclusions are made. The report therefore describes the implemented architecture and defines measurable evaluation procedures, while leaving observed performance values to be completed from actual experiments.
	
	\pagebreak
	%\begin{center}
	%	\Huge \textbf{Abstract} \\
	%\end{center}
	%\par
	
	
	%\pagebreak
	\renewcommand{\contentsname}{\textbf{Table Of Contents} }
	\renewcommand{\cfttoctitlefont}{\huge}
	\tableofcontents\thispagestyle{fancy}
	\setlength{\parindent}{4em}
	\setlength{\parskip}{1em}
	\renewcommand{\baselinestretch}{2.0}
	
	\paragraph{List Of Figures}
	\begin{itemize}
		\item Figure 1: Conceptual resource mismatch motivating IdleGrid.
		\item Figure 2: IdleGrid high-level system architecture.
		\item Figure 3: Job assignment and execution workflow.
		\item Figure 4: Dashboard screen overview (image to be supplied).
		\item Figure 5: Distributed training workflow across worker nodes.
		\item Figure 6: Performance evaluation chart (populate with measured data).
	\end{itemize}
	\paragraph{Acronyms}
	\begin{itemize}
		\item API \hspace{5em}Application Programming Interface
		\item DDP \hspace{5em}Distributed Data Parallel
		\item Gloo \hspace{5em}Collective communications backend for PyTorch
		\item HTTP \hspace{5em}Hypertext Transfer Protocol
		\item LAN \hspace{5em}Local Area Network
		\item REST \hspace{5em}Representational State Transfer
		\item WSL2 \hspace{5em}Windows Subsystem for Linux 2
		
	\end{itemize}
		
	\pagebreak
	\section{Introduction (1 to 1.5 Pages)}
    
Campus laboratories contain computers that are unused at some times or have capacity that is not fully utilized. At the same time, students may have personal devices that are inadequate for compiling large projects, running experiments, processing datasets, or training small machine-learning models. IdleGrid explores whether a campus can make some of its existing computer capacity available through a centrally coordinated, resource-limited service.

The challenge is not simply a shortage of processors. Resources are distributed across multiple PCs, each with a separate operating system, current user, network connection, and local workload. A system that makes the PCs available as one pool must observe the nodes, decide whether a request can fit, tell a worker what to run, bound that work, and present reliable status to the requester. It must also account for the fact that a PC can disappear from the network or become active while a job is running. These characteristics make the project a small distributed-systems problem with both software and operational constraints.

IdleGrid has three principal parts: a Spring Boot Master, a Java Agent installed on each participating Windows PC, and a React web dashboard. Agents periodically report node capacity and pull assigned jobs from the Master. The Master queues ordinary jobs and uses a first-fit policy to select an online node with sufficient reported CPU and memory. Each Agent runs its assigned workload in a Docker container with CPU and memory limits. The dashboard provides node monitoring and screens for jobs, bookings, shared storage, and distributed training. A separate public landing page distributes the portable Windows Agent package; it is not the internal dashboard.

The Master is a control plane rather than a worker. It receives browser requests, records node and job state, schedules eligible ordinary jobs, and makes assignments available to Agents. The Agent is the execution component: it reports capacity, polls for work, launches and watches containers, captures output, and informs the Master about completion. This division allows the central service to coordinate multiple computers without requiring the dashboard to connect to every worker directly. The implementation uses a request/poll pattern, so Agents initiate their own assignment requests.

The implementation also includes a shared-file workflow and a separate CPU-based PyTorch DistributedDataParallel workflow. For shared files, the Master retains a primary copy and queues up to two worker replicas when eligible disk capacity is available. For distributed training, the Master reserves two or three workers and runs one process rank per worker; the user's training code is responsible for data partitioning and gradient synchronization. These features are prototype implementations and have operational dependencies, particularly Windows/WSL2/Docker configuration and direct network connectivity between training workers.

This project should be understood as a prototype rather than a production cloud. Its node, job, booking, and training-run control records are held in memory, and therefore are lost when the Master restarts. File content and the file index are stored on the Master filesystem. Authentication, user identity, quotas, durable job records, checkpointing, replica repair, and a fairness-aware scheduler are outside the current implementation. These boundaries matter when interpreting both the system design and any results.

This report describes only behavior represented in the current repository and README documentation. It distinguishes implemented behavior from planned improvements, and it does not claim performance results that have not been measured. The first chapter introduces the problem, motivation, and significance; the following chapters review related work, identify the gap, specify stakeholders and objectives, describe the system design and algorithms, and present a reproducible performance-evaluation plan. Screenshots, logs, measured values, and institutional particulars are intentionally left for the project team to supply.

	\subsection{Problem Statement (2 - 3 Lines)}

Campus computer laboratories are shared environments rather than dedicated datacenters. Their machines have finite CPU, memory, and disk resources, fluctuate in availability, and must remain usable by their owners. A compute-pooling system must therefore discover participating PCs, keep their status current, match requests to capacity, limit the impact of running workloads, and handle node or network failure. These requirements make resource aggregation more than simply launching a command on a remote computer.

In the current prototype, a central Master maintains node and job records in memory. Agents send heartbeats, and a scheduler assigns a queued job to the first online node whose reported free CPU and memory meet the request. An Agent pulls the assignment and starts a resource-capped container. This centralized first-fit method is intentionally simpler than fairness-aware or datacenter-scale schedulers; it is appropriate for demonstrating a small trusted lab prototype but does not guarantee fairness or globally optimal placement. Figure~\ref{fig:problem-context} illustrates the resource mismatch addressed by the project.

The system has to make decisions from periodically reported data. A heartbeat gives the Master a snapshot, not a guarantee that the reported free resources will remain available until a container starts. Users may submit requests that exceed the capacity of every currently online worker; these should remain queued rather than being silently reduced. A worker can lose its Master connection, a container can exit unsuccessfully, or a DDP peer can be unreachable even when the Agent can contact the Master. These are normal failure conditions in a distributed environment and need to be visible to the user.

There is also a machine-ownership constraint. A lab PC may be used for teaching or administrative work while a background computation is running. Container CPU and memory caps limit the resources allocated to an individual container, but they do not prove that the owner will experience no slowdown. IdleGrid additionally detects user activity and can evict a running job on a transition from idle to active. That is a coarse protection mechanism: it ends the container and reports failure, and it does not preserve progress. The scheduler does not currently enforce an idle-only eligibility rule based on the Agent's activity state.

\begin{figure}[htbp]
\centering
\safeincludegraphics[width=0.88\linewidth,height=6cm,keepaspectratio]{img1.png}
\caption{Conceptual resource mismatch motivating IdleGrid.}
\label{fig:problem-context}
\end{figure}

\textbf{Problem statement:} Design and prototype a campus-LAN platform that discovers participating Windows PCs, monitors their available resources, assigns user-submitted compute jobs to eligible nodes, and executes those jobs within per-container CPU and memory limits, while supporting shared files and a small multi-node training workflow. The system must expose its operation through a web dashboard and accurately report failure and scope limitations. The principal challenges are dynamic node availability, imperfect capacity estimates, non-interference with PC owners, network reachability, and limited persistence and fault recovery.

	
	\subsection{Motivation}

Students who do not own high-performance computers can face barriers when completing compute-intensive coursework. A campus lab already contains machines with useful CPU and memory capacity, but that capacity is not necessarily accessible outside scheduled lab use. A controlled pool could make those resources easier to discover and use without purchasing a separate compute cluster for every student. The potential benefit is reuse of existing equipment, not a claim that the lab has unlimited or cost-free capacity: electricity, administration, maintenance, bandwidth, and staff time remain real costs.

The project is relevant to students, teaching staff, and lab administrators. Students need a simple way to view availability and submit jobs without learning a separate remote-execution interface for every PC. Teaching staff may benefit from a repeatable environment for demonstrations and small experiments, particularly where a lab exercise involves distributed execution. Administrators need visibility into participating nodes, bounded workloads, and clear operational requirements before approving any installation. The public Agent download page can simplify distribution, but it does not replace installation review, endpoint policy, or administrator approval.

IdleGrid is also a practical learning exercise in distributed systems: it combines heartbeats, scheduling, container execution, file replication, web APIs, failure detection, and distributed training. Each feature provides an opportunity to study a distinct systems concern. Heartbeats make liveness observable; polling limits the need for the Master to initiate worker connections; first-fit gives a simple placement baseline; containers apply per-workload limits; and the DDP extension demonstrates that placement alone is not enough when workers must communicate directly.

The prototype is not a production cloud service. In particular, it currently has no authentication or authorization, and therefore must remain on a trusted network and should not be exposed publicly. The dashboard and API do not establish user ownership of jobs or files, and anyone possessing a file key can download the corresponding file. These are important design limitations, not details that can be ignored in a campus deployment. Its value at this stage is demonstrating an end-to-end design and enabling controlled measurement on a small campus deployment after administrative approval.

    \subsection{Significance of the problem}
	
IdleGrid uses a centralized control plane and lightweight worker Agents. The Master is the single point through which the dashboard and Agents coordinate. This design provides a clear source of node and job status and simplifies a small prototype compared with a decentralized membership and scheduling protocol. Its first-fit scheduler is easy to inspect and demonstrate, although it is not intended to provide fairness, optimal packing, or production-scale throughput.

The implementation combines ordinary batch jobs with shared-file handling and a distinct PyTorch DDP workflow. Docker CPU and memory limits bound each ordinary container, while Agent activity detection and job eviction provide a basic response when a user returns after an idle period. This behavior is not a guarantee of zero user impact: the Master currently places jobs using online status and reported capacity, and an eviction is a failure rather than a checkpoint-and-resume operation. The project is therefore best described as a resource-capped lab-compute prototype with activity-aware worker behavior, not as a fully elastic co-tenant scheduler.

The significance of the work is the integration of these capabilities into a single campus-oriented prototype and the opportunity to measure them on real lab equipment. The design is composed of a browser interface, central coordination service, worker-side execution agents, and container runtime, so the project can be discussed through both architectural and empirical perspectives. Unlike an isolated scheduler exercise, the implementation includes complete request paths from dashboard submission to worker execution and returned status, together with file-transfer and DDP workflows.

The principal engineering contribution is integration for a constrained environment: Windows PCs using Docker Desktop and WSL2, a campus LAN, a small node count, and a single central Master. The use of a simple first-fit scheduler keeps placement behavior understandable and supports a baseline implementation against which a future fairness-aware strategy could be evaluated. The use of Docker resource limits gives a concrete mechanism for restricting per-job CPU and memory, while the Agent's activity detector supplies a separate return-to-use response. These mechanisms address different aspects of resource management and should not be conflated: one limits container allocation, while the other can terminate a job after renewed user activity.

The project should be evaluated by measuring rather than assuming its benefits. Evaluation should vary the number of online workers, job CPU and memory requests, workload duration, and worker availability. Relevant metrics include queue-to-assignment delay, job makespan, successful completion ratio, throughput, resource utilization, storage replica completion, and DDP run/artifact success. Owner impact should be measured through a controlled comparison of responsiveness with and without a capped workload, using the same machine and comparable foreground tasks. Repeated trials are required because startup time, image availability, network transfer, background system load, and thermal conditions can affect measurements.

No numerical results are included here because the project documentation does not provide recorded measurements. Accordingly, the report does not claim an improvement over cloud services, existing schedulers, idle-only systems, or a baseline implementation. The Results and Analysis section contains a test matrix and data-collection structure that must be completed from actual runs. Once measurements are collected, they can establish whether the prototype meets its local functional goals and under which test conditions its resource-pooling approach is useful.

	\pagebreak
    
	\section{Literature Review , Gap Identified - 1 to 2 Pages}
The relevant background includes idle-workstation harvesting, volunteer computing, cluster scheduling, container orchestration, and multi-node machine learning. These areas offer useful design patterns, but their assumptions differ from a small Windows-based campus lab.

Condor established the idea of matching computation to idle workstations and is a foundational example of cycle harvesting \cite{condor}. BOINC extended volunteer computing to large pools of independently operated computers, emphasizing client-side participation and fault tolerance \cite{boinc}. These systems demonstrate that distributed spare capacity can be useful, but a campus deployment also needs straightforward local administration, current resource visibility, and a clear way to run containerized coursework workloads.

Cluster schedulers address placement and resource allocation at a larger scale. Dominant Resource Fairness provides a principled way to allocate multiple resource types fairly \cite{drf}, while the Borg experience describes production scheduling and container management at datacenter scale \cite{borg}. IdleGrid currently uses a simpler first-fit policy rather than these more advanced strategies. This keeps placement behavior understandable for a small prototype, but makes fairness and allocation efficiency important future evaluation and development topics.

Container-based co-location work, including Bubble-Up and Heracles, studies the interference that occurs when workloads share hardware and the mechanisms used to control it \cite{bubbleup,heracles}. IdleGrid applies per-container CPU and memory limits and records Windows idle/active transitions, but its implementation does not include feedback-driven interference control. An Agent may evict a running job after a return-to-active transition, and the job is reported as failed rather than resumed. Thus the current design is a practical baseline for experimentation, not a demonstration that user experience is unaffected.

PyTorch DDP supplies a data-parallel training model in which separate processes synchronize gradients; a distributed sampler can assign different training examples to each rank \cite{pytorchddp}. IdleGrid's training workflow orchestrates the placement and launch of one worker process per selected PC, but the submitted program must implement the DDP logic. The system distributes the same uploaded bundle to workers; it does not split data files or automatically convert arbitrary Python scripts into distributed programs.

The systems and concepts reviewed above differ in target scale and assumptions. Condor and BOINC demonstrate resource harvesting and volunteer participation, but their history and operating models do not directly prescribe how a department should deploy a Windows Agent, configure Docker Desktop, and expose a browser dashboard on a private campus LAN. Borg and DRF address larger shared clusters and multi-resource allocation; they provide useful comparison points but would add complexity beyond the current project's first-fit scope. Co-location studies focus on measuring and controlling interference, which makes owner responsiveness an important experiment rather than something that can be inferred from resource caps alone.

\begin{table}[h!]
\centering
\caption{Comparison of related approaches with the IdleGrid prototype.}
\label{tab:related-work}
\begin{tabular}{|p{0.16\linewidth}|p{0.26\linewidth}|p{0.24\linewidth}|p{0.24\linewidth}|}
\hline
Approach & Main focus & Relevance to this project & Difference from current IdleGrid \\
\hline
Condor \cite{condor} & Harvesting cycles from idle workstations & Establishes workstation resource harvesting as a systems problem & IdleGrid is a small campus-LAN prototype with dashboard-based submission and Docker jobs \\
\hline
BOINC \cite{boinc} & Public-resource and volunteer computing & Demonstrates large-scale participation and distributed task execution & IdleGrid assumes a managed lab environment and central Master; it is not a public volunteer network \\
\hline
DRF \cite{drf} & Fair allocation of multiple resource types & Provides a candidate future scheduler and evaluation baseline & Current ordinary placement uses first-fit and does not provide DRF fairness guarantees \\
\hline
Borg \cite{borg} & Large-scale cluster orchestration & Provides context for resource reservation and container management & IdleGrid targets a handful of Windows PCs, not datacenter-scale workloads \\
\hline
Bubble-Up / Heracles \cite{bubbleup,heracles} & Co-location interference and resource efficiency & Motivates controlled measurement of worker-owner impact & IdleGrid has limits and activity-based eviction, but no feedback controller or proven zero-impact guarantee \\
\hline
PyTorch DDP \cite{pytorchddp} & Synchronous data-parallel model training & Supplies the training primitive used by the multi-worker demonstration & IdleGrid coordinates workers; the user script implements model, sampler, and synchronization logic \\
\hline
\end{tabular}
\end{table}

\textbf{Identified gap:} the referenced systems provide important foundations, but do not by themselves document an integrated, small-scale campus prototype combining a web dashboard, Windows Agents, Docker/WSL2 execution, first-fit placement, shared file replicas, and an optional multi-node CPU training workflow. IdleGrid's contribution is this integration and its concrete evaluation opportunity. It is not claimed to outperform Condor, BOINC, or datacenter schedulers; the gap is addressed through a bounded implementation whose limitations can be tested and reported transparently.

\textbf{Positioning and research questions.} The project can be evaluated through questions that match the actual implementation: (1) Can a node heartbeat and job request pass through the full Master--Agent--container lifecycle? (2) How do queue delay and makespan change as the worker count and resource request vary? (3) Under what conditions does a file replica task complete on eligible worker disks? (4) Can the DDP example complete across two or three reachable PCs and produce the requested artifact? (5) What observable change in foreground responsiveness occurs when a capped job runs, and what happens when the Agent detects a return from idle? These questions avoid claiming novelty in the first-fit algorithm itself.
    \pagebreak
    \section{Customer Details and Use Cases}
The primary customers are students who need temporary compute resources for coursework and small experiments but do not have suitable local hardware. Secondary stakeholders are instructors who demonstrate distributed or containerized workloads and laboratory administrators responsible for approving software and network access. The current prototype assumes a trusted campus LAN and administrator approval for installing Java, Docker Desktop, WSL2, and the Agent.

The stakeholder roles have different requirements. A student needs understandable resource choices, clear job status, and access to output without knowledge of the worker's local configuration. An instructor needs a reproducible demonstration and a way to determine whether a job actually completed. A laboratory administrator needs to understand installation privileges, ports, Docker/WSL2 configuration, resource ceilings, and the effect of worker shutdown. A developer/operator needs logs and predictable APIs for diagnosing discovery, scheduling, and execution problems. The current prototype offers only partial support for these operational needs; for example, it provides logs and status screens but not user-specific access control or an administrator policy console.

\textbf{Use case 1 -- Submit a batch job:} A student opens the dashboard, selects a command and CPU/RAM request, and submits it. The Master queues and attempts to place the job. The selected Agent polls for an assignment, starts a Docker container, captures its output, and reports completion. The dashboard displays job status and final output; it is not a live interactive terminal.

\textbf{Use case 2 -- Store and retrieve a file:} A student uploads a file in the Storage page. The Master stores the primary copy and returns a random access key. Eligible online Agents may store up to two replicas. A user with the key can download the file.

\textbf{Use case 3 -- Run distributed training:} A user uploads a ZIP containing a DDP-aware Python entry point and data, selects two or three workers, and submits a command. The Master reserves eligible workers and schedules one rank per node. The program performs sample partitioning and gradient synchronization; rank zero uploads the requested model artifact.

\textbf{Use case 4 -- Monitor lab capacity:} A lab administrator or instructor views node status and reported CPU, memory, and disk capacity in the dashboard and checks whether Agents are reaching the Master.

\textbf{Use case 5 -- Request a targeted session:} A user selects an online node and requests a time-limited SSH-enabled container session through the Book Resources workflow. The backend creates a target-node job and constructs the session command; the requested duration is capped at 24 hours. This is a containerized workload path, not SSH access to the Windows host. The booking records exposed through the booking API are in-memory records and do not implement a durable calendar or reservation guarantee.

\textbf{Use case 6 -- Recover from unavailable capacity:} If no online worker meets the CPU/RAM request, an ordinary job remains queued. The user can inspect its status while the scheduler retries on later passes. The prototype does not guarantee a wait time or fairness among users. If a node becomes stale, the Master marks it offline after the heartbeat timeout and applies the implemented requeue behavior for affected ordinary jobs; DDP worker loss fails the coordinated run.

\textbf{Use-case boundary:} This report describes authorized prototype use on a trusted campus network. It does not specify anonymous public job submission, multi-tenant identity, billing, quotas, arbitrary network exposure, or guaranteed long-running services. Such use cases would require security and operations features not present in the implementation.
	\pagebreak
	\section{Project Objectives}
    \begin{enumerate}
      \item To design and implement a centralized platform for monitoring compute resources on campus laboratory PCs.
      \item To enable submission, first-fit scheduling, and Docker-based execution of CPU- and memory-constrained jobs.
      \item To provide a shared-file workflow with a Master-stored primary copy and best-effort worker replicas.
      \item To implement a dashboard workflow for coordinating CPU-based PyTorch DDP runs across two or three workers.
      \item To define and conduct controlled measurements of job execution, resource use, storage replication, and worker availability.
    \end{enumerate}
	\pagebreak	
	\section{Proposed Method/Algorithm }
The problem is to accept a compute request and place it on an eligible campus PC while keeping the worker, job, and user-visible status coordinated. A solution must account for delayed or missing heartbeats, resource requests that exceed all reported capacity, Agent polling intervals, and the difference between ordinary batch work and coordinated multi-node training. IdleGrid addresses these constraints with a central Master, periodic heartbeats, a first-fit scheduler for ordinary jobs, pull-based Agent assignments, and separate gang placement for DDP runs. The current solution favors simplicity and traceability over fairness, durability, and optimal placement.

	\subsection{Problem Definition}

Given a set of worker nodes \(N\), each with a reported online state and free CPU and memory, and a user job request \(J=(c_j,m_j,\mathrm{command},\mathrm{image})\), the Master should select an eligible node \(n\in N\) whose reported available CPU and memory satisfy \(c_n\ge c_j\) and \(m_n\ge m_j\), or leave the job queued if no such node exists. The expected outputs are an updated job state, an assignment for the selected Agent to pull, and a completion result reported after container execution. For a DDP request, the inputs include a validated training bundle, worker count, resource requests, container image, and output filename; the expected result is coordinated execution on the selected workers and, on success, a downloadable rank-zero artifact.

The ordinary-workload input can include an optional target node and an optional duration for a time-limited SSH-style container session. A request also depends on the worker having access to the selected container image and a usable workspace. The platform does not accept a generic high-level program and infer its parallelism: for distributed training, the uploaded ZIP must contain a compatible Python entry point and the command supplied by the user must be valid in the training container. The run request specifies a worker count of two or three and a requested amount of CPU and memory for each selected worker.

The output is not limited to whether the command started. The Master exposes the current state of a job or run; the Agent returns the container result and captured output when ordinary execution completes; and the DDP page provides per-rank status and the final artifact when rank zero successfully uploads it. For storage, the output of an upload is metadata and an access key, while replica acknowledgements are tracked separately from the durable Master copy. This distinction is important: an upload completing to the Master does not by itself prove that all worker replicas have been created.

The objective is to demonstrate end-to-end resource monitoring, placement, execution, and status reporting on a small campus LAN. Constraints include Windows worker PCs, Java 17+, Docker Desktop with WSL2, network access to the Master, and direct inter-worker TCP reachability for DDP. The prototype's node/job/booking/run state is in memory; capacity reports may become stale; first-fit is not a fairness policy; ordinary job output is available after completion rather than as a live terminal; and the system has no authentication. These assumptions make it unsuitable for public or untrusted deployment without further engineering.

	\subsection{Proposed Idea/System}
The core idea is to turn existing lab PCs into visible worker nodes without making the browser communicate directly with each PC. The dashboard submits requests to the Master, the Master tracks worker heartbeats and job state, and each Agent pulls its own assignment. Docker enforces per-job CPU and memory limits at container launch. This separation simplifies the control path and allows Agents to operate behind the central API.

The first-fit scheduler walks online nodes and assigns a queued job to the first node with sufficient reported CPU and memory after accounting for active reservations. Heartbeats older than 15 seconds cause a node to be marked offline. Jobs that are affected by node loss may be requeued according to the Master workflow; no checkpoint is created. For shared files, the Master stores the authoritative file and metadata and queues replica tasks for eligible workers. DDP placement is a separate operation: two or three eligible workers are reserved together and each runs a rank using host networking. The novelty claimed here is the integrated, lab-focused prototype and its workflow, not a new scheduling algorithm.
    
	\subsection{System Architecture}
IdleGrid uses a centralized architecture. The React dashboard calls the Spring Boot Master over HTTP. The Master maintains in-memory node, job, booking, and training-run state, stores primary files and an index on its filesystem, and runs scheduling and placement logic. Java Agents on Windows PCs report heartbeats and pull jobs or storage tasks. Assigned jobs run in Docker Desktop containers through WSL2. For DDP, selected worker containers also need direct peer-to-peer TCP connectivity. Figure~\ref{fig:architecture} shows these boundaries; the implementation description follows the project technical guide \cite{idlegridguide}. The centralized Master and separate Agents are the principal integration choices of this prototype.

\textbf{Dashboard.} The internal dashboard is a React application built with Vite and Tailwind CSS. Its pages include an overview, job submission and status, a terminal-style batch command interface, resource booking, shared storage, and distributed training. Browser requests go to the Master API; the dashboard is not an execution agent and does not directly invoke Docker on a worker. In development, the Vite server proxies API requests to the configured backend. The public download site is a separate application and serves the portable Agent ZIP.

\textbf{Master service.} The Master is implemented in Java with Spring Boot. Controllers expose node, job, booking, storage, and distributed-training endpoints. The scheduler evaluates ordinary queued jobs and reserves a group of workers for DDP. In-memory stores hold node/job/booking/run control data. File bytes and an index are stored under the configured filesystem directory, while training bundles and artifacts are stored under the configured training directory. The Master also advertises its LAN endpoint through the configured rendezvous/discovery mechanisms. The Master is a single coordination point; its failure interrupts new API operations and its in-memory records are not recovered on restart.

\textbf{Worker Agent.} A Java background process runs on each participating Windows PC. Scheduled tasks send heartbeats, check Windows input inactivity, poll for job assignments, monitor running containers, and request storage replication tasks. The Agent keeps a stable node identifier in a local file unless explicitly configured otherwise. It launches jobs using Docker CLI commands, captures session output in local logs, and reports running/completed status to the Master. The Agent is packaged as a portable ZIP with a launcher, configuration, and setup helper; Java 17+ and Docker Desktop/WSL2 are runtime requirements.

\textbf{Container runtime and data paths.} Ordinary jobs run in Docker containers with CPU and memory limits and a per-job workspace mounted at \texttt{/workspace}. Ordinary containers use Docker's default bridge network and may publish image-exposed ports. DDP containers request host networking because rank processes must communicate directly. The worker runtime is not a native Windows container environment: Docker Desktop uses WSL2. This deployment makes the platform dependent on installed versions, Docker configuration, firewall policy, and campus network topology.

\textbf{Storage component.} The Master maintains a canonical file copy and JSON metadata index. After an upload, it may queue up to two replica tasks on online Agents with available disk. Agents poll, download the source bytes from the Master, write local copies, and acknowledge completion. This is not a distributed object store: the repository does not implement MinIO, S3, erasure coding, user buckets, or a repair/rebalance service. DDP bundles and output artifacts follow their own training-run paths and have separate validation and size constraints.

\textbf{External communication and ports.} The HTTP Master API uses TCP port 9090 by default. Local Agent discovery uses UDP port 9091. The default DDP rendezvous uses TCP port 29500, with additional peer connections depending on runtime behavior. Agents must reach the Master; multi-node DDP also requires direct worker-to-worker connectivity. A successful heartbeat therefore confirms only one network path and does not prove that the DDP communication path is available.

\begin{figure}[htbp]
\centering
\safeincludegraphics[width=0.92\linewidth,height=7cm,keepaspectratio]{img2.png}
\caption{High-level architecture of the IdleGrid prototype.}
\label{fig:architecture}
\end{figure}

	\subsection{System Analysis and Design}
The design follows the repository's implemented API boundaries and operating assumptions. The subsections below list functional and non-functional requirements, explain the main workflow, define test cases for execution, and summarize the placement algorithm. Tests are described as planned/expected checks; their results must be recorded during actual runs.
	%\pagebreak
	\subsubsection{Requirement Specification}
\textbf{Functional requirements:}
\begin{enumerate}
  \item The Master shall accept Agent heartbeats and expose node status and reported capacity to the dashboard.
  \item A user shall be able to submit a job with a command and CPU/memory request and inspect its state and completion output.
  \item An Agent shall pull an assignment, start a resource-limited Docker container, and report job completion or failure.
  \item A user shall be able to upload and download shared files using a generated access key; eligible worker nodes may store replicas.
  \item A user shall be able to submit and inspect a DDP run on two or three workers and retrieve its artifact when produced.
\end{enumerate}
\textbf{Non-functional requirements and constraints:} Java 17 or later is required for the backend and Agent; worker PCs require Windows, Docker Desktop, and WSL2. The dashboard is browser-based. Agents must reach the Master on the campus network; DDP additionally requires worker-to-worker TCP connectivity and compatible Docker host networking. Capacity values are periodic estimates. The prototype has no authentication, durable control database, or checkpoint-based restart and should be restricted to a trusted network.

Table~\ref{tab:requirements} makes the requirements traceable to observable behavior. Priorities describe the prototype's intended functional importance, not a contractual production service level.
\begin{table}[h!]
\centering
\caption{Functional and operational requirement traceability.}
\label{tab:requirements}
\begin{tabular}{|p{0.1\linewidth}|p{0.43\linewidth}|p{0.24\linewidth}|p{0.12\linewidth}|}
\hline
ID & Requirement & Verification evidence & Priority \\
\hline
FR1 & Register/update nodes from Agent heartbeats and expire stale presence & Node summary and heartbeat test & High \\
\hline
FR2 & Accept a job with CPU, memory, command, optional image and target & Job API response and dashboard & High \\
\hline
FR3 & Place a job only when an online node meets reported CPU/RAM & Assignment status and scheduler logs & High \\
\hline
FR4 & Execute an assignment in a bounded Docker container & Agent log and container inspection & High \\
\hline
FR5 & Return status and completion output to the Master & Job status endpoint and dashboard & High \\
\hline
FR6 & Store file content on Master and expose key-based retrieval & Upload/download checksum comparison & Medium \\
\hline
FR7 & Schedule eligible worker replicas and track acknowledgements & Storage page and Agent logs & Medium \\
\hline
FR8 & Reserve 2--3 workers for a DDP run and expose rank status/artifact & Run status and rank-zero artifact & Medium \\
\hline
NFR1 & Operate within approved Windows/WSL2/Docker campus environment & Deployment record & High \\
\hline
NFR2 & Keep user-facing pages separate from per-node execution & Architecture and API trace & High \\
\hline
NFR3 & Make failure and prototype limitations explicit to operators & Logs, status and report caveats & High \\
\hline
NFR4 & Restrict prototype access to a trusted network until auth exists & Deployment review & Critical \\
\hline
\end{tabular}
\end{table}

	\subsubsection{Flowcharts / DFDs / UML / Functional Block Diagram}
The ordinary job data flow is: user submits a request to the dashboard; the dashboard sends it to the Master; the Master queues and attempts first-fit placement; the selected Agent polls and receives the assignment; Docker runs the container; and the Agent sends completion status and output to the Master for dashboard display. If no node meets the request, the job remains queued. Figure~\ref{fig:job-flow} is reserved for a workflow illustration.

The heartbeat path runs independently of job submission. At its configured interval, each Agent samples local information and sends a heartbeat containing its node identity, IP address, CPU and memory totals/free values, disk information, and status. The Master updates the node record and last-seen time. A stale heartbeat eventually causes the node to be considered offline. This mechanism provides liveness information but is not a continuous resource reservation: the host may change resource use between measurements.

The storage path begins when the browser uploads a multipart file to the Master. The Master stores the file and associates it with a randomly generated key in the JSON index. Where disk capacity permits, the storage service queues replica work. An Agent pulls a replica task, fetches the primary content, writes it to its configured storage directory, and sends a completion acknowledgement. The primary Master copy remains the authoritative source. If an Agent is offline or does not have enough disk, a worker replica may not be created.

The DDP path differs from an ordinary job because a run requires a group of workers to start coordinated ranks. The Master validates the uploaded archive and request, finds the requested number of eligible PCs, reserves them together, and creates a rank assignment for each. Each Agent runs the same training bundle in the configured CPU worker image. The user's script initializes the Gloo process group, partitions examples (for example with \texttt{DistributedSampler}), and wraps the model using DDP. Rank zero writes the requested model artifact and uploads it to the Master. Failures in a rank or the peer network can fail the group; there is no checkpoint-based recovery.

\textbf{Logical data-flow boundary.} The browser exchanges HTTP requests and responses with the Master. Agents exchange heartbeats, assignment polls, job state updates, and storage replica messages with the Master. Worker containers receive assignments through their Agents and, for DDP only, exchange distributed-training traffic directly with peer workers. The user's browser does not open a direct execution channel to a worker. Keeping these paths distinct aids diagnosis: a successful dashboard API call, a live Agent heartbeat, a successful ordinary job, and a successful multi-node DDP run each validate different parts of the system.

\begin{figure}[htbp]
\centering
\safeincludegraphics[width=0.88\linewidth,height=6.5cm,keepaspectratio]{img3.png}
\caption{Ordinary job submission, placement, and execution workflow.}
\label{fig:job-flow}
\end{figure}

	\subsubsection{Design and Test Cases}
Table~\ref{tab:testcases} lists functional checks that should be executed in a controlled environment. The expected behavior is derived from the documented API and workflows; the observation column must be completed from actual tests. Before each test, record the software versions, node identifiers, network configuration, and whether images are already cached. This avoids mistaking an image-pull delay or a network policy issue for scheduler performance.
\begin{table}[h!]
\centering
\caption{Functional test cases and observations to be recorded.}
\label{tab:testcases}
\begin{tabular}{|p{0.08\linewidth}|p{0.25\linewidth}|p{0.35\linewidth}|p{0.2\linewidth}|}
\hline
ID & Scenario & Expected result & Observation \\
\hline
TC1 & Agent heartbeat & Node appears online with reported capacity & [Run and record] \\
\hline
TC2 & Submit eligible job & Job progresses through assignment and completion & [Run and record] \\
\hline
TC3 & Insufficient capacity & Job remains queued until eligible capacity is available & [Run and record] \\
\hline
TC4 & Upload and retrieve file & Key is returned and content can be downloaded & [Run and record] \\
\hline
TC5 & DDP run on 2--3 workers & Worker ranks report status and rank zero artifact is downloadable & [Run and record] \\
\hline
TC6 & Stop worker heartbeats & Node expires after timeout; affected work follows documented failure/requeue behavior & [Run and record] \\
\hline
TC7 & Restart Master & In-memory control records are absent after restart; local file store persists when configured path remains intact & [Run and record] \\
\hline
TC8 & Request a targeted node & Scheduler considers only the selected target where requested & [Run and record] \\
\hline
TC9 & Return from idle during a job & Agent behavior and job failure report match configured activity/eviction behavior & [Run and record] \\
\hline
TC10 & Upload file with worker replica capacity & Master stores primary and Agent reports successful replica placement & [Run and record] \\
\hline
TC11 & DDP worker network unavailable & Run failure is visible; no successful artifact is claimed & [Run and record] \\
\hline
TC12 & Repeated identical job trials & All trial outcomes and elapsed times are recorded, including failures & [Run and record] \\
\hline
\end{tabular}
\end{table}

The functional tests should be followed by controlled performance trials rather than one-off demonstrations. Run each workload repeatedly with a fixed input, image, resource request, and worker configuration. Separate cold-start trials (where an image or application needs initialization) from warm trials (where the image is already present). Do not combine them into one average unless the test explicitly represents a mixed deployment. Save job IDs, run IDs, timestamps, Agent logs, Master logs, and the final output or artifact for each trial.

Where possible, use a small deterministic command to validate the lifecycle before using a longer CPU workload. Confirm the result content, not only the status label. For file transfer, compare a cryptographic checksum of the source file with the downloaded file. For DDP, record every rank's status and inspect the artifact; submission of a run alone is not evidence that distributed training completed correctly. For owner-impact evaluation, obtain lab approval and use a harmless repeatable foreground activity, with the same machine and comparable background conditions for baseline and treatment runs.

	\subsubsection{Algorithms and Pseudo Code}
The ordinary job placement procedure is summarized below. It captures the current first-fit policy and does not imply a fairness guarantee.
\begin{verbatim}
for each queued job J:
    for each node N in registry order:
        if N is online and reported free CPU/RAM
           after active reservations meets J's request:
            assign J to N
            stop searching for J
    if no eligible node was found:
        leave J queued
\end{verbatim}
Agents independently poll for assignments, launch the configured Docker image with the requested CPU and memory limits, monitor the container, and report the result. DDP uses a separate group-placement path that reserves the requested workers together before launching one rank per worker.

The scheduler's node eligibility predicate can be stated as: the node is online; any target-node constraint is satisfied; reported free CPU and memory less active reservations are at least the request; and the node is not otherwise excluded by the current scheduling logic. The figures used in this predicate are observations rather than hard host-level reservations. Docker applies container limits after assignment, but does not make the scheduler's earlier capacity report atomic with the host's actual load.

\textbf{Ordinary-job state model.} A job begins in \texttt{QUEUED}. When a suitable node is selected, it becomes \texttt{ASSIGNED} and records the worker. When the Agent starts the container and reports the transition, it becomes \texttt{RUNNING}. A successful completion is \texttt{DONE}; an unsuccessful container exit, watchdog failure, or activity-triggered eviction is reported as \texttt{FAILED}. Offline-node handling can return affected ordinary work to the queue under the implemented Master behavior. The exact observed transition should be verified from logs and API responses because process/network timing can affect when a transition is visible.

\textbf{Complexity and design trade-off.} If a scheduling pass examines \(J\) queued jobs and scans up to \(N\) registered nodes for each job, the upper-bound comparison work is proportional to \(J\times N\). For the intended small demonstration pool this is straightforward and inexpensive; there is no benchmark claim for larger \(N\). First-fit may leave some capacity fragmented or may favor nodes earlier in the registry iteration order. Sorting by resource efficiency, fair shares, load prediction, queue aging, or data locality would alter the policy and should be introduced only with an explicit design and comparative evaluation.

\textbf{DDP group placement.} The DDP workflow must find a complete set of two or three eligible workers before starting the ranks. This is a gang-scheduling constraint: a partial group is not sufficient for a coordinated run. The same bundle is made available to each selected rank; data partitioning belongs to the submitted Python program. At the application level, DDP synchronizes gradients; it does not merge independently trained model files after the fact. Run completion therefore depends on compatible software, rank initialization, network paths, and rank-zero artifact upload.

\textbf{Data and state design.} A node record represents a worker known to the Master and includes a stable node ID, IP address, reported CPU/RAM/disk capacity, status, and last-heartbeat time. A job record contains its generated ID, requested CPU and RAM, command/image information, status, optional assigned node, and result. Booking records contain a node, requester, interval, and requested resources, but are in-memory and do not enforce a persistent reservation calendar. Distributed-training run state includes run-level configuration and per-worker/rank status. File metadata links a generated key and file properties to the Master-stored content and any known replica nodes. This separation reflects the different durability mechanisms: in-memory control state versus filesystem-backed shared file content.

\textbf{API interaction design.} Dashboard and Agent clients use the Master as the API boundary. Node summary and job/run status are read operations. Heartbeat, job submission, state updates, file upload, replica acknowledgement, and DDP submission are write or transfer operations. Assignment retrieval uses Agent polling; a no-work response is an expected condition rather than a failure. The detailed endpoint list and typical request shape are included in Appendix~B. The endpoint contract does not add authentication: deployment must treat API access as trusted-network access.

\textbf{Configuration and installation design.} A Master may be discovered through the configured rendezvous service, UDP discovery, or the local-subnet fallback; an explicit Master URL is available when automatic discovery cannot identify the correct interface. The portable Agent package keeps the launcher, executable JAR, and properties file together. The setup helper prepares Windows features and Docker Desktop, but requires elevated privileges and can require a reboot. Therefore, installation is an operational procedure requiring lab administrator approval and verification on a non-critical test machine before wider deployment.

    \subsubsection{UI/UX Wireframes where applicable}
The dashboard includes overview, booking, jobs, storage, and distributed-training pages. The supplied screenshot should be captured from the running application and should show the actual implemented UI rather than a generated mockup. Figure~\ref{fig:dashboard} is a placeholder for that evidence.

The overview page is intended to answer whether workers are online and what resource capacity they report. The job and terminal-style pages let a user submit commands and inspect their status and output after completion; the word "terminal" describes a batch command bridge and does not mean an interactive, streaming shell. The booking page allows a user to target an online node and request a time-limited container session, while the separate booking API stores simple in-memory booking records. The storage page shows file metadata and access keys and supports upload/download. The distributed-training page accepts a ZIP bundle and configuration, and displays the run and rank status and artifact link. Each screenshot should be labelled with the page name and should avoid exposing access keys, private network addresses, or personal data.

The UX should present states explicitly: online/offline for nodes and queued/assigned/running/done/failed for jobs. For no-capacity conditions, a queued state should not be presented as if execution has started. For storage, the availability of the Master primary and the count/status of worker replicas should remain distinct. For DDP, users should be informed that all selected workers need a compatible image and peer network connectivity, and that the uploaded script must implement DDP itself. These distinctions help avoid overpromising behavior not present in the implementation.

\begin{figure}[htbp]
\centering
\safeincludegraphics[width=\textwidth,height=0.62\textheight,keepaspectratio]{img4.png}
\caption{Illustrative IdleGrid dashboard view with five simulated worker nodes and sample resource availability.}
\label{fig:dashboard}
\end{figure}
	%\subsubsection{sequence diagram}
	\pagebreak
	\section{Performance Study}
	\subsection{Implementation/Simulation Environment}

The documented implementation uses Java 17 and Spring Boot 3.3.4 for the Master, a Java 17 background Agent packaged for Windows, and React 18, Vite, and Tailwind CSS for the dashboard. Worker jobs run in Docker Desktop using its WSL2 backend. The separate distributed-training image provides a CPU-based PyTorch/Gloo runtime. The Master exposes HTTP on port 9090 and local Agent discovery uses UDP port 9091. The dashboard development server uses port 5173. These technologies are already integrated in the project; the report does not propose a replacement stack.

The deployment targeted by the project context is a small prototype of one Master and up to four worker PCs. Ordinary job tests can vary the number of online workers from one to four and the CPU/RAM request, subject to the capacity of the actual lab machines. DDP currently supports two or three workers. Record the exact PC models, operating-system versions, CPU and RAM, Docker/WSL2 versions, network type, and software versions in the experiment log before collecting data. The following are suggested test conditions, not claimed experimental results:
\begin{itemize}
  \item Worker count for ordinary jobs: 1, 2, 3, and 4 online worker nodes where available.
  \item Job resource requests: record at least two CPU limits and two memory limits that every selected worker can support; a small smoke-test request may use 1 CPU and 512 MB RAM.
  \item Job type: short deterministic commands for scheduling checks, followed by a representative CPU workload for timing and resource observation.
  \item Node availability: all selected workers online, then one worker stopped during a controlled run to observe failure handling.
  \item DDP worker count: 2 and 3, using the same bundle, image, training command, and dataset for comparable trials.
\end{itemize}

Table~\ref{tab:environment} provides a per-machine inventory to fill in before the experiment. Hardware variation matters because the workers may not be identical. A difference in processor generation, memory pressure, SSD/HDD speed, or network connection can affect both scheduling and execution time. For reproducibility, give each node a stable test label such as Master, Worker 1, and Worker 2 in the report and retain the actual node ID only in private logs if it should not be published.
\begin{table}[h!]
\centering
\caption{Experimental environment inventory (complete from the actual test machines).}
\label{tab:environment}
\begin{tabular}{|p{0.17\linewidth}|p{0.22\linewidth}|p{0.22\linewidth}|p{0.22\linewidth}|}
\hline
Property & Master PC & Worker 1 & Worker 2 / additional workers \\
\hline
PC/model identifier & [Record] & [Record] & [Record] \\
\hline
Processor / logical CPUs & [Record] & [Record] & [Record] \\
\hline
Installed RAM & [Record] & [Record] & [Record] \\
\hline
Operating system & [Record] & [Record] & [Record] \\
\hline
Storage type / free space & [Record] & [Record] & [Record] \\
\hline
Docker Desktop / WSL2 version & [Record] & [Record] & [Record] \\
\hline
Network connection & [Record] & [Record] & [Record] \\
\hline
\end{tabular}
\end{table}

Define the following measurements for each completed trial. Let \(T_s\) be the time the Master accepts a job submission, \(T_a\) the time it assigns a worker, \(T_r\) the time the container is reported running, and \(T_f\) the time completion is reported. Queue-to-assignment delay is \(T_a-T_s\). Assignment-to-start delay is \(T_r-T_a\). Container makespan is \(T_f-T_r\); end-to-end job latency is \(T_f-T_s\). Separating these values distinguishes scheduling and polling overhead from workload execution, provided the timestamps are captured consistently.

For a fixed set of submissions, completion ratio is \(N_{\mathrm{successful}}/N_{\mathrm{submitted}}\), where a successful job must reach the documented success state and produce its expected output. Throughput is \(N_{\mathrm{successful}}/T_{\mathrm{window}}\), where the observation window is explicitly defined. Mean CPU utilization is the time-average of measured worker CPU use divided by the worker's logical CPU capacity; memory utilization is the time-average of used memory divided by total memory. Record the sampling interval and source of these utilization measurements. Storage replica completion ratio is completed replicas divided by replica tasks issued during the trial, with the primary upload outcome reported separately. DDP speedup is \(T_1/T_p\) only when a comparable single-worker baseline \(T_1\) and a \(p\)-worker runtime \(T_p\) use the same model, dataset, software image, epoch count, and batch-size policy. A failed or timed-out run is not silently discarded; report it as a failure and state how it was handled in summary statistics.

Owner impact should use an explicitly defined task instead of a subjective claim alone. Possible observations include foreground task response time, completion time of a fixed interactive operation, and a short user-rated smoothness scale collected with consent. Record the workload's requested CPU/RAM, whether the node was active or idle at start, whether eviction occurred, and the foreground baseline. The report should avoid presenting a single informal observation as proof of general non-interference.

	\subsection{Results and Analysis}
No measured performance dataset or test-run log was provided with the README files. Numerical results, charts, and claims of speedup or improvement are therefore intentionally left blank. Execute the test cases on the target lab PCs and complete Table~\ref{tab:results} with observed values before submitting the report. Record failed trials as well as successful ones, and state environmental conditions that could explain variation.

Use the following sequence for each experiment group. First, verify all selected Agents are online and record their reported capacity. Second, submit one fixed workload and capture the returned job/run identifier and server-side timestamps or available log timestamps. Third, follow it through assignment, container start, and completion, saving the output and relevant logs. Fourth, verify correctness by checking the expected command output, file checksum, or DDP artifact. Finally, repeat the same configuration enough times to characterize variability and document every failure. If the implementation does not expose a timestamp needed for a metric, state the measurement method and instrument consistently rather than implying precision the system does not provide.

For ordinary scheduling, test a workload that is sufficiently long to separate execution time from startup overhead, then run shorter jobs to characterize overhead. Keep the requested CPU and memory values constant when comparing worker counts. When comparing resource requests, keep the same node set and workload. Use separate tables or charts for cold-image and warm-image runs because pulling a container image adds network and startup cost. When measuring throughput, define whether jobs are submitted serially or concurrently and report the queue length; otherwise throughput values from different experiments are not directly comparable.

For shared storage, use known file sizes and verify downloaded bytes with a checksum. Record upload latency, download latency, the number of replica tasks created, successful acknowledgements, and the number of workers eligible at the start. The implementation's replica target is up to two worker copies, not a guarantee of two copies for every file. For DDP, use the same data bundle and training parameters for all worker counts, record the image digest or precise image version if available, and confirm the artifact file can be downloaded and opened. A DDP run that completes with two workers establishes a functional test for that setup, but not a general speedup claim.

\begin{table}[h!]
\centering
\caption{Performance observations to complete from actual experiments; no values are prefilled.}
\label{tab:results}
\begin{tabular}{|p{0.25\linewidth}|p{0.22\linewidth}|p{0.22\linewidth}|p{0.2\linewidth}|}
\hline
Measurement & Test configuration & Observed value & Notes / trial count \\
\hline
Queue-to-assignment delay & [Workers, request] & [Measure] & [Record] \\
\hline
Job makespan & [Workload, request] & [Measure] & [Record] \\
\hline
Job completion ratio & [Submitted jobs] & [Measure] & [Record] \\
\hline
Worker CPU/memory utilization & [Worker and workload] & [Measure] & [Record] \\
\hline
Storage replica completion ratio & [Files and online nodes] & [Measure] & [Record] \\
\hline
DDP runtime / speedup & [2 or 3 workers] & [Measure] & [Record] \\
\hline
Owner responsiveness impact & [Idle/active controlled trial] & [Measure] & [Record] \\
\hline
\end{tabular}
\end{table}

For each row in Table~\ref{tab:results}, attach a trial sheet containing: trial number; start/end date and local time; Master and worker labels; node count; CPU and RAM request; image version and cache state; workload command and input checksum; job or run ID; measured timestamps; final state; output/artifact verification; and relevant error/log summary. Report central tendency and variability when multiple trials are available. For small samples, show the individual observations as well as any mean or median so the reader can see outliers.

Figure~\ref{fig:ddp-flow} is a placeholder for an explanatory workflow image; Figure~\ref{fig:results-chart} is reserved for a chart produced from the measured data in Table~\ref{tab:results}. Do not use generated or illustrative data as experimental evidence.

\begin{figure}[htbp]
\centering
\safeincludegraphics[width=0.88\linewidth,height=6.5cm,keepaspectratio]{img5.png}
\caption{Distributed training workflow implemented by the prototype.}
\label{fig:ddp-flow}
\end{figure}

\begin{figure}[htbp]
\centering
\fbox{\parbox[c][4cm][c]{0.88\linewidth}{\centering\textbf{Performance chart to be added after experiments}}}
\caption{Measured performance comparison (to be generated from recorded experimental data).}
\label{fig:results-chart}
\end{figure}

The interpretation of these measurements should separate scheduler overhead from workload execution, and should report whether any apparent improvement is repeatable. A higher worker count does not necessarily reduce runtime for a small workload because startup, bundle transfer, communication, and network overhead may dominate. Any owner-impact conclusion must be based on a controlled observation rather than inferred from the presence of Docker limits.

The analysis should answer the question represented by each experiment rather than merely restating a table. If assignment delay rises with worker count, inspect Master scheduling and heartbeat processing; if container makespan changes, consider CPU contention, memory pressure, input/output behavior, and hardware differences. If a DDP run with more workers is slower, discuss the work size, synchronization overhead, image startup, and network transfer before attributing the result to an algorithm. If replicas are not created, distinguish insufficient worker disk from offline workers, delayed polling, or transfer failure. Discuss both expected behavior and surprises, and include unsuccessful runs as operational evidence.

\textbf{Threats to validity.} A small number of lab machines cannot establish datacenter-scale behavior. Heterogeneous hardware and background user activity can confound comparisons. Image pulls and training-bundle uploads can dominate short-run times. Clock differences across PCs can distort timestamps unless time synchronization is verified; prefer Master-side timestamps for cross-node lifecycle measures where possible. Network congestion, Wi-Fi isolation, firewall rules, and Docker host networking affect DDP independently of compute capacity. Finally, a synthetic image-classification bundle validates the workflow but does not establish model quality on a real application dataset.

	\subsection{Summary of performance study}

The current report establishes the metrics and experiment matrix but contains no measured values. It cannot yet determine whether first-fit placement, DDP, or resource capping performs better under any tested condition. The ordinary scheduler performs a linear scan over registered nodes for each queued job, giving \(O(JN)\) placement work for \(J\) queued jobs and \(N\) nodes in a pass; this is suitable for a small prototype but has not been benchmarked at larger scale. Complete the result table and analyze actual observations before drawing conclusions about throughput, scalability, or user impact.
    
    \subsection{Limitations}
The prototype targets a small trusted campus LAN and Windows worker PCs using Docker Desktop and WSL2; results may not generalize to other operating systems, networks, or large clusters. Node, job, booking, and training-run control state is held in memory and is lost when the Master restarts. There is no authentication or authorization, and a file access key grants download access to its holder. Worker file replicas are best-effort and are not backed by a repair or rebalance guarantee.

Reported capacity can become stale between heartbeat and scheduling, and the first-fit scheduler does not provide fairness or optimal packing. The Agent's activity detection and return-to-active eviction do not guarantee negligible owner impact; an evicted job is failed rather than checkpointed and resumed. Ordinary job output is collected after completion, not streamed through a live interactive terminal. DDP depends on compatible images, Docker Desktop host networking, and direct inter-worker TCP connectivity; loss of a DDP worker can fail the run. The report's performance conclusions remain incomplete until repeatable measurements are collected on approved lab hardware.

\pagebreak
\section{Conclusions and Future Work}

IdleGrid implements a campus-focused prototype for discovering worker PCs, monitoring reported resources, placing ordinary jobs with a first-fit scheduler, running those jobs in Docker containers, managing shared files with best-effort worker replicas, and coordinating CPU-based PyTorch DDP runs across a small worker group. These features address the engineering objective of integrating compute monitoring, execution, storage, and distributed-training workflows behind one Master and dashboard. However, the supplied project documentation contains no measured experiment results, so the effectiveness, performance, and owner impact cannot yet be judged quantitatively. The prototype's known constraints include in-memory control state, absence of authentication, approximate capacity reports, limited fault recovery, and network/configuration dependencies. Consequently, the system is a working-scope prototype whose performance and suitability for lab use must be confirmed through controlled tests and administrative review.

Future work should first collect repeatable measurements on the target PCs and improve validation and observability. Further engineering could add authentication, per-user ownership and quotas, persistent control data, stronger failure recovery, checkpoint/restart support, replica repair, and fairness-aware scheduling such as DRF. DDP reliability could be improved through automated network checks and clearer recovery behavior. Any expansion beyond the campus prototype should also evaluate isolation, privacy, resource contention, and administrator approval requirements.

\pagebreak

\section{Scholarly Contributions}
Students should list the relevant scholarly contributions arising from the project under the following categories:
The README documentation supplied for this draft does not identify any publications, filed patents, exhibitions, or certifications. Complete each item only with verified project outputs; do not treat planned or unpublished work as an achieved contribution.
\begin{enumerate}
	\item Research papers (published, accepted, communicated).
	\item Patents or copyright filed.
	\item Project exhibitions.
	\item Open-source contributions, GitHub repository, demo video, deployed product live link.
	\item Certifications relevant to the project.
\end{enumerate}


	\pagebreak
	
	\section{References}
	\begin{thebibliography}{9}
\bibitem{condor}
M. Litzkow, M. Livny, and M. Mutka,
\textit{Condor: A Hunter of Idle Workstations},
Proceedings of the 8th International Conference on Distributed Computing Systems, 1988.

\bibitem{boinc}
D. P. Anderson,
\textit{BOINC: A System for Public-Resource Computing and Storage},
Proceedings of the 5th IEEE/ACM International Workshop on Grid Computing, pp. 4--10, 2004.

\bibitem{drf}
A. Ghodsi, M. Zaharia, B. Hindman, A. Konwinski, S. Shenker, and I. Stoica,
\textit{Dominant Resource Fairness: Fair Allocation of Multiple Resource Types},
Proceedings of the 8th USENIX Symposium on Networked Systems Design and Implementation, 2011.

\bibitem{bubbleup}
J. Mars, L. Tang, R. Hundt, S. Skadron, and M. L. Soffa,
\textit{Bubble-Up: Increasing Utilization in Modern Warehouse Scale Computers via Sensible Co-locations},
Proceedings of the 44th Annual IEEE/ACM International Symposium on Microarchitecture, 2011.

\bibitem{heracles}
A. Lo, L. Cheng, R. Govindaraju, P. Ranganathan, and C. Kozyrakis,
\textit{Heracles: Improving Resource Efficiency at Scale},
Proceedings of the 42nd Annual International Symposium on Computer Architecture, 2015.

\bibitem{borg}
B. Burns, B. Grant, D. Oppenheimer, E. Brewer, and J. Wilkes,
\textit{Borg, Omega, and Kubernetes},
Communications of the ACM, vol. 59, no. 5, 2016.

\bibitem{pytorchddp}
PyTorch Contributors,
\textit{DistributedDataParallel},
PyTorch Documentation, available at \url{https://pytorch.org/docs/stable/generated/torch.nn.parallel.DistributedDataParallel.html}.

\bibitem{idlegridguide}
IdleGrid project team,
\textit{IdleGrid -- Project and Technical Guide},
Project repository documentation, DETAILED\_README.md.
\end{thebibliography} 

\pagebreak
\appendix
\section{Appendix A: Plagiarism Report, AI Similarity Report}
This appendix contains the plagiarism and AI similarity report of the project. 

\pagebreak
\section{Appendix B: API and Data-Flow Reference}
This appendix summarizes the implemented Master API surfaces described in the project technical guide. It is included to support system understanding and testing; exact field validation and response structures should be checked against the running version of the repository.

\begin{table}[h!]
\centering
\caption{Node and ordinary-job API summary.}
\label{tab:api-jobs}
\begin{tabular}{|p{0.22\linewidth}|p{0.28\linewidth}|p{0.36\linewidth}|}
\hline
Method and path & Caller & Purpose \\
\hline
\texttt{POST /nodes/heartbeat} & Agent & Create or refresh node capacity and last-seen status \\
\hline
\texttt{GET /nodes/summary} & Dashboard & List known node records and reported status/capacity \\
\hline
\texttt{POST /jobs/submit} & Dashboard/test client & Create a job and immediately attempt scheduling \\
\hline
\texttt{GET /jobs} & Dashboard & List jobs held by the current Master process \\
\hline
\texttt{GET /jobs/\{id\}/status} & Dashboard & Read one job's state, assignment, and result \\
\hline
\texttt{POST /jobs/\{id\}/running} & Agent & Report container start and related runtime information \\
\hline
\texttt{POST /jobs/\{id\}/complete} & Agent & Report result and success/failure outcome \\
\hline
\texttt{GET /agent/\{nodeId\}/assignment} & Agent & Pull an assignment; no assignment is returned as an empty response \\
\hline
\end{tabular}
\end{table}

\begin{table}[h!]
\centering
\caption{Storage and distributed-training API summary.}
\label{tab:api-storage}
\begin{tabular}{|p{0.22\linewidth}|p{0.28\linewidth}|p{0.36\linewidth}|}
\hline
Method and path & Caller & Purpose \\
\hline
\texttt{GET /storage/summary} & Dashboard & Summarize Master/worker disk capacity and file count \\
\hline
\texttt{GET /storage/files} & Dashboard & List file metadata and access keys \\
\hline
\texttt{POST /storage/files} & Dashboard & Upload multipart file to Master storage \\
\hline
\texttt{GET /storage/files/\{key\}/content} & Browser/user & Download Master-stored content by key \\
\hline
\texttt{GET /agent/\{nodeId\}/storage/assignment} & Agent & Pull a file replica task \\
\hline
\texttt{GET /agent/storage/\{key\}/content} & Agent & Download source bytes for a replica \\
\hline
\texttt{POST /agent/\{nodeId\}/storage/\{key\}/complete} & Agent & Acknowledge replica placement result \\
\hline
\texttt{POST /distributed-training/runs} & Dashboard & Submit a validated bundle and worker configuration \\
\hline
\texttt{GET /distributed-training/runs/\{runId\}} & Dashboard & Read run and per-worker/rank status \\
\hline
\texttt{GET /distributed-training/runs/\{runId\}/bundle} & Worker container & Download the submitted training bundle \\
\hline
\texttt{POST /distributed-training/runs/\{runId\}/artifact} & Rank zero & Upload the final model artifact \\
\hline
\texttt{GET /distributed-training/runs/\{runId\}/artifact} & Dashboard/user & Download the completed artifact \\
\hline
\end{tabular}
\end{table}

The Master HTTP API uses port 9090 by default. Local Agent discovery uses UDP port 9091, and the default DDP rendezvous port is TCP 29500. The DDP worker may need additional direct connections according to the distributed runtime. Firewall and campus network rules must be verified by the system administrator; a successful heartbeat does not establish that peer-to-peer DDP traffic is permitted.

\pagebreak
\section{Appendix C: Deployment and Experiment Record}
The following checklist is intended for an authorized lab setup and a repeatable demonstration. It is not a substitute for institutional approval, security review, or examination of the setup script before running it with administrator privileges.

\textbf{Pre-deployment checklist:}
\begin{enumerate}
  \item Obtain approval to install and run the Agent on each selected PC and to run Docker Desktop/WSL2.
  \item Select one Master and label the worker nodes; record hardware, OS, Java, Docker Desktop, and network information in Table~\ref{tab:environment}.
  \item Start the Master and verify the backend is listening on the configured TCP port. Confirm its advertised LAN URL is reachable from each worker.
  \item Install Java 17+ and Docker Desktop with WSL2 on each worker. Run the provided setup helper only after reviewing its effects and coordinating any required restart.
  \item Configure automatic or manual Master discovery, start each Agent, and verify heartbeat visibility in the dashboard and Master logs.
  \item Pull or locally build the exact ordinary-job and DDP image tags required by the test. Do not assume a public registry package contains a tag until an image pull has been confirmed.
  \item For DDP, verify host networking and peer TCP reachability among the selected workers before scheduling a training run.
  \item Use non-sensitive, synthetic test data; keep the API restricted to the approved campus network because the prototype has no authentication.
\end{enumerate}

\textbf{Per-trial record (copy one record for each run):}
\begin{table}[h!]
\centering
\caption{Blank experimental trial record.}
\label{tab:trial-record}
\begin{tabular}{|p{0.35\linewidth}|p{0.5\linewidth}|}
\hline
Field & Value to record \\
\hline
Trial ID / date / operator & [Complete] \\
\hline
Master and worker labels & [Complete] \\
\hline
Software and image versions & [Complete] \\
\hline
Workload / command / input checksum & [Complete] \\
\hline
Requested CPU / memory / worker count & [Complete] \\
\hline
Image cache state / network notes & [Complete] \\
\hline
Job ID or DDP run ID & [Complete] \\
\hline
Submission / assignment / start / finish times & [Complete] \\
\hline
Final status and output/artifact verification & [Complete] \\
\hline
Errors, logs, exclusions, or deviations & [Complete] \\
\hline
\end{tabular}
\end{table}

\textbf{Evidence to include in the final submitted copy:} a genuine screenshot of the running dashboard; a redacted node summary; a successful job's state progression and output; a storage upload/download verification; a DDP run status and artifact verification if the multi-node test is available; and the completed trial tables and plots. Replace every bracketed field with verified information or explicitly mark the test as not performed. Do not leave illustrative prompts or placeholder boxes in the final printed report.

\end{document}
